\chapter{Fundamentos del controlador PID}
\label{ch:fundamentos-pid}

\section{Introducción del capítulo}
El controlador proporcional--integral--derivativo constituye una de las herramientas más
utilizadas para regular variables de proceso en ingeniería. Su permanencia se explica por la
relación favorable entre sencillez estructural, interpretabilidad física y capacidad de
adaptación a problemas reales. En sistemas eléctricos aislados, donde la frecuencia debe
recuperarse con rapidez y sin oscilaciones excesivas, comprender el papel de cada acción
del PID resulta imprescindible antes de discutir cualquier método de sintonización.

\section{Objetivo del capítulo}
Este capítulo presenta los principios de funcionamiento del PID y describe cómo la acción
proporcional, integral y derivativa afectan la respuesta transitoria y estacionaria del sistema.
El propósito es establecer una base conceptual sólida para que, en los capítulos posteriores,
la optimización de las ganancias no se interprete como un ajuste ciego, sino como una
decisión informada sobre velocidad, error residual y amortiguamiento.

\section{Estructura general del controlador PID}
En su forma más conocida, el PID genera una señal de control a partir del error entre la
referencia y la salida medida. La acción proporcional responde al error instantáneo; la
integral acumula el historial del error para corregir desviaciones persistentes; y la derivativa
anticipa la tendencia futura al observar la rapidez con que cambia la señal. Esta combinación
permite construir un regulador flexible, capaz de equilibrar rapidez y estabilidad cuando sus
parámetros se eligen de forma adecuada.

\section{Acción proporcional}
La ganancia proporcional $K_p$ establece la intensidad con la que el controlador reacciona
ante un error presente. Cuando $K_p$ aumenta, el sistema suele responder con mayor rapidez,
pero también crece el riesgo de sobreimpulso y de oscilaciones, sobre todo en plantas con
retardos, dinámica no mínima o acoplamientos hidráulicos. En una microred aislada, un
$K_p$ demasiado bajo vuelve lenta la recuperación de frecuencia, mientras que uno
excesivamente alto puede provocar ciclos de corrección agresivos sobre el gobernador.

\section{Acción integral}
La acción integral, ponderada por $K_i$, tiene como misión eliminar el error en estado
estacionario. Su efecto es especialmente valioso cuando existen perturbaciones sostenidas o
sesgos estructurales que impiden alcanzar exactamente la referencia. No obstante, una
integración demasiado fuerte puede introducir oscilaciones de baja frecuencia y fenómenos de
windup cuando el actuador trabaja cerca de sus límites físicos, por lo que su ajuste debe
considerar tanto precisión como restricciones de operación.

\section{Acción derivativa}
La acción derivativa, asociada a $K_d$, actúa como un mecanismo de amortiguamiento porque
responde a la velocidad de cambio del error. En términos prácticos, aporta un efecto de
``frenado anticipatorio'' que ayuda a reducir rebotes y a suavizar la trayectoria de retorno a la
referencia. Su utilidad es notoria en sistemas donde la planta tiende a comportarse de forma
subamortiguada; sin embargo, también es sensible al ruido de medición, lo que obliga a usarla
con criterio, especialmente en implementaciones digitales.

\section{Formulación matemática}
En tiempo continuo, la ley de control puede escribirse como
\[
u(t)=K_p e(t)+K_i \int_0^t e(\tau)\,d\tau +K_d \frac{de(t)}{dt}.
\]
En el dominio de Laplace, esta expresión conduce a la función de transferencia del
controlador
\[
G_c(s)=K_p+\frac{K_i}{s}+K_d s.
\]
Cuando el controlador se implementa en un PLC o en una plataforma digital, la acción integral
y derivativa se aproximan mediante diferencias finitas y tiempos de muestreo definidos. Esa
versión discreta es la que normalmente se usa en simulaciones y prototipos aplicados a
microredes modernas.

\section{Ventajas y limitaciones}
Entre las principales fortalezas del PID destacan su bajo costo computacional, su facilidad de
implementación y la posibilidad de interpretar físicamente cada parámetro. Por ello, sigue
siendo el estándar industrial incluso en aplicaciones complejas. Su limitación principal no es
la estructura en sí, sino la dependencia crítica del método de sintonización: cuando la planta
es fuertemente no lineal o trabaja bajo perturbaciones variables, un ajuste deficiente degrada
rápidamente el desempeño y puede anular las ventajas del controlador.

\section{Cierre del capítulo}
Comprender el PID como una combinación de rapidez, memoria y amortiguamiento permite
entender por qué su ajuste no puede reducirse a una fórmula universal. El siguiente capítulo
profundiza precisamente en los criterios con que se juzga si una sintonización es buena o
deficiente, es decir, en las métricas que permiten comparar objetivamente distintas
configuraciones del controlador.
